\section{Benchmark problem definition}\label{sec:problem}

All quantities are in SI units. Wherever possible the definitions of
\citet{Arostica2025} (hereafter AR25) are reused verbatim so that the solid
sub-problem is identical to their Benchmark~1.

% ===========================================================================
\subsection{Geometry}\label{sec:geometry}

\paragraph{Myocardium $\Os$} The truncated ellipsoid of AR25 (their eqs.~9--11),
with long axis along $x$ and apex at $x<0$:
\begin{equation}
\vect{x}(\mu,\theta) = \bigl(r_{\mathrm{long}}\cos\mu,\;
 r_{\mathrm{short}}\sin\mu\cos\theta,\;
 r_{\mathrm{short}}\sin\mu\sin\theta\bigr),\qquad \theta\in[-\pi,\pi],
\end{equation}
with the dimensions of \cref{tab:geometry}. The base $\Gbase$ joins the
endocardial and epicardial rims by a straight line in each meridian plane.

\paragraph{Fluid domain $\Of$} The LV cavity enclosed by $\Gendo$ and closed at
the endocardial rim by the planar cap $\Gcap$ ($x = x_{\mathrm{cap}}$). The cap
contains an inflow (mitral) port $\Gin$ and an outflow (aortic) port $\Gout$,
both circular and fixed in space; the rest of the cap,
$\Gcap\setminus(\Gin\cup\Gout)$, is a no-slip wall.
\TODO{Confirm port radii/centres and whether short inflow/outflow tubes are
extruded from the cap (recommended: tube length $\geq 2D$ to keep outflow
boundary away from the vortex; state it). Values below are a proposal.}
\TODO{State how the cap-wall/endocardium junction moves (cap rim follows the
endocardial rim via mesh motion; cap remains planar).}

\begin{table}[t]
\centering
\caption{Geometry. Myocardial dimensions are those of AR25; derived values are
for the unloaded reference configuration.}
\label{tab:geometry}
\begin{tabular}{lll}
\toprule
Quantity & Symbol & Value \\
\midrule
Endocardium, short / long semi-axis & $r_{\mathrm{short}}^{\mathrm{endo}}, r_{\mathrm{long}}^{\mathrm{endo}}$ & \SI{25}{mm}, \SI{90}{mm} \\
Endocardium, $\mu$ range & & $[-\pi,\,-\arccos(5/17)]$ \\
Epicardium, short / long semi-axis & $r_{\mathrm{short}}^{\mathrm{epi}}, r_{\mathrm{long}}^{\mathrm{epi}}$ & \SI{35}{mm}, \SI{97}{mm} \\
Epicardium, $\mu$ range & & $[-\pi,\,-\arccos(5/20)]$ \\
Base plane, endocardial rim & $x_{\mathrm{cap}}$, $r_{\mathrm{cap}}$ & \SI{26.47}{mm}, \SI{23.89}{mm} \\
Base plane, epicardial rim & & \SI{24.25}{mm}, \SI{33.89}{mm} \\
Cavity volume (reference) & $V_0$ & \SI{168.3}{mL} \\
Wall thickness, equator / apex & & \SI{10}{mm} / \SI{7}{mm} \\
\midrule
Inflow port radius / centre & $R_{\mathrm{in}}$, $\vect{c}_{\mathrm{in}}$ & \TODO{\SI{11}{mm}, $(x_{\mathrm{cap}}, +11.5, 0)$\,mm} \\
Outflow port radius / centre & $R_{\mathrm{out}}$, $\vect{c}_{\mathrm{out}}$ & \TODO{\SI{9}{mm}, $(x_{\mathrm{cap}}, -13, 0)$\,mm} \\
\bottomrule
\end{tabular}
\end{table}

\begin{figure}[t]
\centering
\fbox{\parbox[c][5cm][c]{0.9\linewidth}{\centering
\TODO{Fig.~1: (a) myocardium with $\Gendo,\Gepi,\Gbase$, points $\vect{p}_0,\vect{p}_1$;
(b) fluid domain with $\Gcap,\Gin,\Gout$; (c) fibre and sheet fields;
(d) cut through meshes M1--M4.}}}
\caption{Benchmark geometry, boundaries and meshes.}
\label{fig:geometry}
\end{figure}

% ===========================================================================
\subsection{Fibre architecture}
Fibres $\f$ and sheets $\s$ follow AR25 Section~3.2 exactly: a transmural
coordinate $\bar t$ from $\Delta\bar t=0$ ($\bar t=0$ on $\Gendo$, $1$ on
$\Gepi$, zero flux on $\Gbase$), an ellipsoidal basis with interpolated radii,
and a helix angle $\alpha(\bar t) = -60^\circ + 120^\circ\,\bar t$.

% ===========================================================================
\subsection{Solid model}\label{sec:solid}
In the reference configuration $\Os$, the displacement $\uu$ satisfies
\begin{equation}
\rho_{\mathrm{s}}\,\ddot{\uu} - \Div\PP = \vect 0 \quad\text{in } \Os,
\qquad \PP = \FF\SSS,
\end{equation}
with the visco-active Holzapfel--Ogden second Piola--Kirchhoff stress of AR25
(their eqs.~2--6):
\begin{equation}
\SSS = \frac{\partial\Psi_{\mathrm{aniso}}}{\partial\EE}
     + \eta\,\dot{\EE} + \tau(t)\,\f\otimes\f,
\end{equation}
where $\Psi_{\mathrm{aniso}}$ includes the smooth fibre-compression switch
($k=100$) and the volumetric penalty $\tfrac{\kappa}{4}(J^2-1-2\ln J)$, and the
active stress $\tau(t)$ solves AR25 eq.~(5). Parameters are listed in
\cref{app:parameters}.

\TODO{Decide the timing of $\tau(t)$ for the FSI cycle. Option A: keep AR25
($t_{\mathrm{sys}}=\SI{0.16}{s}$, $t_{\mathrm{dias}}=\SI{0.484}{s}$) and put
filling \emph{before} $t=0$ (preload phase). Option B (current xenosim/solids4foam
setup): filling in $[0,t_1]$, activation shifted to $t_{\mathrm{sys}}=t_1$.
Option B is more physiological; state the shift explicitly.}

% ===========================================================================
\subsection{Fluid model}\label{sec:fluid}
Blood is an incompressible Newtonian fluid. In ALE form on the moving domain
$\Of$ with mesh velocity $\ww$:
\begin{align}
\rho_{\mathrm{f}}\Bigl(\left.\frac{\partial\vv}{\partial t}\right|_{\chi}
 + \bigl((\vv-\ww)\cdot\nabla\bigr)\vv\Bigr)
 - \dive\sig_{\mathrm{f}} &= \vect 0, &
\dive\vv &= 0 \quad\text{in } \Of,\\
\sig_{\mathrm{f}} &= -p\,\II + \mu_{\mathrm{f}}\bigl(\nabla\vv + \nabla\vv^{\mathsf T}\bigr),
\end{align}
with $\rho_{\mathrm f}=\SI{1050}{kg/m^3}$ and
$\mu_{\mathrm f}=\SI{3.5e-3}{Pa.s}$ ($\nu_{\mathrm f} = \SI{3.33e-6}{m^2/s}$).
The flow is assumed laminar. \TODO{Report peak Reynolds and Womersley numbers
based on port diameters; justify laminar assumption.}

% ===========================================================================
\subsection{Interface and boundary conditions}\label{sec:bcs}

\paragraph{Fluid--solid interface $\Gfsi=\Gendo$} Kinematic and dynamic
continuity,
\begin{equation}
\vv = \dot{\uu},\qquad
\sig_{\mathrm f}\,\nn = J^{-1}\PP\NN\,\frac{\mathrm dA}{\mathrm da}
\quad\text{on }\Gfsi,
\end{equation}
replacing the prescribed endocardial pressure of AR25.

\paragraph{Solid} On $\Gepi$ and $\Gbase$ the Robin (spring--dashpot)
conditions of AR25 (normal-only on $\Gepi$ with
$\alpha_{\mathrm{epi}}=\SI{1e8}{Pa/m}$, $\beta_{\mathrm{epi}}=\SI{5e3}{Pa.s/m}$;
all directions on $\Gbase$ with $\alpha_{\mathrm{top}}=\SI{1e5}{Pa/m}$,
$\beta_{\mathrm{top}}=\SI{5e3}{Pa.s/m}$).

\paragraph{Fluid} No slip on the cap wall. On the ports, a plug
(or parabolic \TODO{choose}) velocity profile with prescribed volumetric flow
rates $Q_{\mathrm{in}}(t)$ and $Q_{\mathrm{out}}(t)$, both defined positive
(\cref{sec:phases}); a port is a no-slip wall when its flow rate is zero.
The pressure level is fixed by \TODO{reference-pressure strategy: mean cavity
pressure, or pressure at one point}.

% ===========================================================================
\subsection{Cardiac phases without valves}\label{sec:phases}
Valves are replaced by the time-dependence of $Q_{\mathrm{in}}$ and
$Q_{\mathrm{out}}$ (\cref{fig:waveforms}):
\begin{description}
\item[{Phase I --- passive filling, $t\in[0,t_1]$:}]
  $\tau=0$; $Q_{\mathrm{in}}(t) = Q_{\mathrm E}\sin^2\!\bigl(\pi t/t_1\bigr)$
  \TODO{or E-wave + A-wave}, $Q_{\mathrm{out}}=0$.
\item[{Phase II --- isovolumic contraction, $t\in[t_1,t_2]$ (optional):}]
  $Q_{\mathrm{in}}=Q_{\mathrm{out}}=0$; $\tau$ rises.
\item[{Phase III --- ejection, $t\in[t_2,t_3]$:}]
  $Q_{\mathrm{in}}=0$; $Q_{\mathrm{out}}(t) = Q_{\mathrm S}\sin^2\!\bigl(\pi (t-t_2)/(t_3-t_2)\bigr)$.
\item[{Phase IV --- relaxation, $t\in[t_3,T]$:}] $Q_{\mathrm{in}}=Q_{\mathrm{out}}=0$; $\tau$ decays.
\end{description}
The filled and ejected volumes are
$V_{\mathrm{fill}} = Q_{\mathrm E}t_1/2$ and $V_{\mathrm{ej}} = Q_{\mathrm S}(t_3-t_2)/2$.

\TODO{Key design decision, see \cref{sec:disc:dilemma}. With flow rates
prescribed on all open boundaries the cavity volume is kinematically imposed,
$V(t)=V_0+\int_0^t(Q_{\mathrm{in}}-Q_{\mathrm{out}})\,\mathrm ds$, and the
cavity pressure becomes a Lagrange multiplier determined by the solid. This is
a demanding but well-posed test of the coupling (incompressibility dilemma).
Consider adding a variant (``Benchmark 3b'') in which ejection is driven by a
pressure/three-element Windkessel outlet so that $V(t)$ is an output.}

\begin{table}[t]
\centering
\caption{Phase timings and flow-rate parameters. \TODO{fill from the xenosim runs.}}
\label{tab:phases}
\begin{tabular}{llll}
\toprule
Parameter & Symbol & Value & Note \\
\midrule
End of filling & $t_1$ & \TODO{\SI{0.65}{s}} & current tutorial \\
End of isovolumic contraction & $t_2$ & \TODO{} & \\
End of ejection & $t_3$ & \TODO{} & \\
Cycle length & $T$ & \SI{1.0}{s} & \\
Peak filling rate & $Q_{\mathrm E}$ & \TODO{} & \\
Peak ejection rate & $Q_{\mathrm S}$ & \TODO{} & \\
Filled / ejected volume & $V_{\mathrm{fill}}, V_{\mathrm{ej}}$ & \TODO{} & must be equal for a closed cycle\\
Activation onset & $t_{\mathrm{sys}}$ & \TODO{} & AR25 eq.~(6) \\
\bottomrule
\end{tabular}
\end{table}

\begin{figure}[t]
\centering
\fbox{\parbox[c][4cm][c]{0.9\linewidth}{\centering
\TODO{Fig.~2: $Q_{\mathrm{in}}(t)$, $Q_{\mathrm{out}}(t)$, $\tau(t)$ and resulting $V(t)$ over one cycle, phases shaded.}}}
\caption{Driving waveforms over one cardiac cycle.}
\label{fig:waveforms}
\end{figure}

% ===========================================================================
\subsection{Benchmark steps}\label{sec:steps}
\begin{description}
\item[Step 0 --- solid only.] AR25 Benchmark~1, Step~1 (active + passive,
  prescribed endocardial pressure). Verifies the solid solver against nine
  reference codes.
\item[Step 1 --- prescribed-motion CFD.] The fluid problem alone with the
  endocardial motion of Step~0 imposed as a Dirichlet ALE boundary condition
  \TODO{or with an analytical wall motion}. Verifies the moving-domain fluid
  solver and mesh motion.
\item[Step 2 --- coupled passive filling.] Phase~I only, $\tau=0$.
  Exercises added mass and the passive anisotropic solid.
\item[Step 3 --- coupled full cycle.] Phases~I--IV.
\end{description}

% ===========================================================================
\subsection{Quantities of interest}\label{sec:qoi}
\begin{enumerate}
\item cavity volume $V(t)$ (from the fluid mesh and, independently, from
  $\Gendo$) --- checks volume conservation;
\item mean cavity pressure $\bar p(t)$ and the pressure at the apex
  $p_{\mathrm{apex}}(t)$, and the intraventricular pressure difference
  $\Delta p_{\mathrm{IVPD}}(t) = p_{\mathrm{base}}-p_{\mathrm{apex}}$;
\item the pressure--volume loop and stroke work $\oint \bar p\,\mathrm dV$;
\item displacements $\uu(\vect p_0,t)$, $\uu(\vect p_1,t)$ at the AR25 points
  $\vect p_0=(0.025,0.03,0)$, $\vect p_1=(0,0.03,0)$\,m, and of the epicardial apex;
\item fluid kinetic energy $E_k(t) = \tfrac12\rho_{\mathrm f}\int_{\Of}|\vv|^2$
  and viscous dissipation $\Phi(t)$;
\item vortex metrics in the mid-plane $z=0$: circulation, vortex-core position
  \TODO{criterion: $\lambda_2$ or $Q$, threshold};
\item velocity profiles along prescribed lines at given times \TODO{define lines};
\item coupling statistics: FSI iterations per step, wall-clock time, cores.
\end{enumerate}
Discrepancy between solutions is measured with the relative discrepancy
$\RED$ of AR25 (eq.~21), sampled every \SI{10}{ms}.
